GSAS-II, a widely-used package for powder and single-crystal
diffraction analysis, is supported by an extensive set of tutorials, help pages,
and reference documents totalling $\sim$500,000 words across five source categories comprising 331 HTML files.
Navigating this wealth of documentation presents challenges for new and experienced users alike.
We present \textit{Query-GSAS}, a retrieval-augmented generation (RAG)
system that provides natural-language access to the complete GSAS-II
documentation corpus.
Source documents are ingested into a local vector database; at query time, semantically relevant passages are retrieved by dense embedding search and a large language model (LLM) synthesises a fluent, cited answer from those passages.
Four backends are supported: an Ollama server (using any model available via
\texttt{ollama pull}), \texttt{llama-cpp-python} for in-process GGUF
inference requiring no separate server process, an Anthropic API option, and a retrieval-only mode that returns raw documentation passages without LLM synthesis.
A post-processing citation injector connects every substantive claim in
the generated response to the original GSAS-II source document by inserting
an inline marker providing a reference to a ranked bibliography below the response.
In the web UI, these are clickable hyperlinks.
\textit{Query-GSAS} is incorporated into GSAS-II using a separate wxPython dialog window, but can also be run as a separate utility offering alternative user interfaces, including a command-line tool and a browser-based interface that uses a self-contained FastAPI web application. \textit{Query-GSAS} can be run as a fully local architecture, which satisfies the security and data-residency requirements of research facilities that require air-gapped or high-performance computing environments. The package serves as a model for how to develop natural-language documentation search tools for other large scientific packages. 